\section{Interference balances and source response}\label{sec:interference}

The local normal form describes what a collision does. The integral
representation describes what must balance for it to occur. We derive
two complementary descriptions, using momentum damping and spatial
heat averaging. Neither description, by itself, forbids the event.

\subsection{Positive momentum source, signed cosine observation}
At the collision, differentiation of \eqref{eq:F-H} gives
\begin{align}
 \int_0^\infty F_{\lambda_*}(u)\cos(x_*u)\dd u&=0,\label{eq:cos-balance}\\
 \int_0^\infty uF_{\lambda_*}(u)\sin(x_*u)\dd u&=0,\label{eq:sin-balance}\\
 \int_0^\infty u^2F_{\lambda_*}(u)\cos(x_*u)\dd u&=-a\ne0,
 \label{eq:curvature-balance}
\end{align}
where \(\lambda_*=\lambda_0-\tau_*\).
Thus a strictly positive source can produce a double node through two
signed observation balances. The second moment in
\eqref{eq:curvature-balance} measures the transverse crossing.

Fix the observation point \(x_*\), and put
\[
 c_\pm(u)=\max\{\pm\cos(x_*u),0\},\quad
 A_\pm(\tau)=\int_0^\infty F_{\lambda_0-\tau}(u)c_\pm(u)\dd u,
\]
\[
 B_\pm(\tau)=
 \frac{\int_0^\infty u^2F_{\lambda_0-\tau}(u)c_\pm(u)\dd u}
      {A_\pm(\tau)}.
\]
Since \(x_*\ne0\), both cosine signs occupy open intervals and
\(A_\pm>0\). Dominated differentiation in time yields
\begin{equation}\label{eq:sector-damping}
 A_\pm'=-A_\pm B_\pm,\qquad
 \frac{\dd}{\dd\tau}\log\frac{A_+}{A_-}=-(B_+-B_-).
\end{equation}
At the event \(A_+=A_-=:A\), and
\begin{equation}\label{eq:sector-curvature}
 a=-A(B_+-B_-).
\end{equation}
If \(a>0\), the negative cosine sector has the larger mean \(u^2\)
and therefore damps faster. The observed amplitude rises through zero.
If \(a<0\), the positive sector damps faster and the crossing is downward.
Equation~\eqref{eq:sin-balance} is the additional stationarity constraint;
the sector-mass equality alone is only a zero condition.

These means use the \emph{linear} source, not the squared-amplitude
measure \eqref{eq:spectral-measure}. The corners of \(c_\pm\) cause no
difficulty for the time derivatives above: the observation point is
fixed and the source moments dominate the integrands. No spatial
differentiation through moving sector boundaries is being used.

\subsection{A spatial heat-average balance}
Choose \(0<\Delta\leq\tau_*\), and define
\[
 k_\Delta(y)=(4\pi\Delta)^{-1/2}\ee^{-y^2/(4\Delta)},\qquad
 f(y)=h(\tau_*-\Delta,y).
\]
The heat multiplier gives \(h(\tau_*,x)=k_\Delta*f(x)\).
This identity first holds in \(L^2\); the integral and all its displayed
derivatives are continuous by Cauchy--Schwarz against Gaussian
derivatives, so it also holds pointwise. Define a signed measure
\[
 \dd\mu(y)=k_\Delta(x_*-y)f(y)\dd y.
\]
Its absolute moments of every fixed order are finite. Differentiating
the Gaussian proves
\begin{equation}\label{eq:gaussian-moments}
 \int\dd\mu=0,\qquad
 \int(y-x_*)\dd\mu=0,\qquad
 \int(y-x_*)^2\dd\mu=4\Delta^2a.
\end{equation}
Indeed, \(\partial_x k_\Delta(x-y)=(y-x)k_\Delta/(2\Delta)\);
the second derivative is
\(\bigl((y-x)^2/(4\Delta^2)-1/(2\Delta)\bigr)k_\Delta\).
The zeroth moment vanishes, removing the last term.

Let \(\mu=\mu_+-\mu_-\) be the positive and negative parts.
Since \(a\ne0\), neither is zero. The first two identities imply
\[
 \mu_+(\R)=\mu_-(\R)=:M>0,\qquad
 \frac1M\int y\dd\mu_+=\frac1M\int y\dd\mu_-=:m.
\]
Writing \(\sigma_\pm^2=M^{-1}\int(y-m)^2\dd\mu_\pm\), the last
identity becomes
\begin{equation}\label{eq:variance-balance}
 a=\frac{M}{4\Delta^2}(\sigma_+^2-\sigma_-^2).
\end{equation}
Thus the earlier positive and negative amplitudes, after Gaussian
weighting about the future event, have equal masses and centroids.
Their variance difference supplies the nonzero crossing curvature.
This is an exact conditional description of a collision, valid at
any such earlier time; it is not an estimate that excludes one.

\subsection{Smooth displacement under source perturbation}
The divergent velocity of an individual zero does not imply singular
dependence of the \emph{joint event} on source parameters.
Let \(q(u)\) be a real, even, bounded smooth function whose derivatives
have polynomial growth, and perturb the supplied source to
\[
 F_{\lambda,\eta}(u)=\ee^{\lambda u^2+\eta q(u)}\Phi(u).
\]
The moment bounds and Schwartz properties persist on compact
parameter sets. The resulting field \(h_\eta\) is analytic in
\((\tau,x,\eta)\), with
\[
 r(\tau,x):=\left.\partial_\eta h_\eta(\tau,x)\right|_{\eta=0}
 =\int_0^\infty F_{\lambda_0-\tau}(u)q(u)\cos(xu)\dd u.
\]

\begin{proposition}[Event response]\label{prop:response}
A collision of \(h_0\) continues uniquely as an analytic local
event \((\tau_*(\eta),x_*(\eta))\). At \(\eta=0\),
\begin{equation}\label{eq:event-response}
 \tau_*'=-\frac{r}{a},\qquad
 x_*'=\frac{br}{a^2}-\frac{r_x}{a},
\end{equation}
where the right-hand side is evaluated at the unperturbed event.
\end{proposition}
\begin{proof}
Apply the analytic implicit-function theorem to
\((h_\eta,\partial_xh_\eta)=0\). Its Jacobian in \((\tau,x)\) is
\[
 \begin{pmatrix}h_\tau&h_x\\h_{x\tau}&h_{xx}\end{pmatrix}_{*}
 =\begin{pmatrix}a&0\\b&a\end{pmatrix},
\]
with determinant \(a^2\ne0\). Differentiation of the two event equations
gives \(a\tau_*'+r=0\) and \(b\tau_*'+ax_*'+r_x=0\), proving
\eqref{eq:event-response}.
\end{proof}

This is the source-response calculation of
\cite{SourceResponse} in the present time convention.
It is local at a nondegenerate event, not uniform as \(a\to0\).
A general tilt changes the canonical theta arithmetic. The proposition
therefore explains event stability and sensitivity without furnishing
a permitted deformation of the canonical source or proving that its
event exists.
